\documentclass[11pt,a4paper]{article}

% ------------------------------------------------------------
% Pacotes
% ------------------------------------------------------------
\usepackage[utf8]{inputenc}
\usepackage[T1]{fontenc}
\usepackage[brazil]{babel}
\usepackage{lmodern}
\usepackage{microtype}
\usepackage{amsmath,amssymb}
\usepackage{enumitem}
\usepackage{geometry}
\usepackage{titlesec}
\usepackage{booktabs}
\usepackage{tabularx}
\usepackage{hyperref}

% ------------------------------------------------------------
% Layout
% ------------------------------------------------------------
\geometry{
    top=1.8cm,
    bottom=1.8cm,
    left=1.9cm,
    right=1.9cm
}

\setlength{\parindent}{0pt}
\setlength{\parskip}{3pt}

\setlist[itemize]{
    leftmargin=1.35em,
    itemsep=0.8pt,
    topsep=2pt,
    parsep=0pt
}

\setlist[enumerate]{
    leftmargin=1.55em,
    itemsep=0.8pt,
    topsep=2pt,
    parsep=0pt
}

\titleformat{\section}
    {\bfseries\normalsize}
    {}
    {0pt}
    {}

\titlespacing*{\section}
    {0pt}
    {7pt}
    {2pt}

\titleformat{\subsection}
    {\bfseries\itshape\normalsize}
    {}
    {0pt}
    {}

\titlespacing*{\subsection}
    {0pt}
    {5pt}
    {1pt}

\hypersetup{
    colorlinks=true,
    linkcolor=black,
    urlcolor=blue
}

% ------------------------------------------------------------
% Documento
% ------------------------------------------------------------
\begin{document}

\begin{center}
    {\large\bfseries MO438 -- Redes Complexas}\\[2pt]
    \textbf{Proposta de Projeto (versão revisada)}
\end{center}

\vspace{-2mm}

% ============================================================
\section*{1. Título}

Da rede lexical à rede contextual: reorganização das representações
internas de Transformers com Redes Complexas

% ============================================================
\section*{2. Breve descrição do projeto}

Transformers representam tokens por vetores em espaços de alta dimensão. Na
entrada do modelo, todas as ocorrências de um mesmo tipo de token possuem o
mesmo \textit{embedding} lexical. À medida que atravessam as camadas do
Transformer, essas representações incorporam informações do contexto e podem se
reorganizar no espaço vetorial.

Este projeto propõe representar essa transformação como uma sequência de
redes complexas. Os vértices corresponderão a tipos de token ou a ocorrências
de tokens (Seção~5), e as arestas conectarão representações semelhantes. A
comparação das redes permitirá investigar em que medida a estrutura permanece
organizada pela identidade lexical e em que medida passa a refletir o
contexto ao longo das camadas.

A hipótese central não é que a identidade lexical desapareça, mas que exista
uma reorganização progressiva das vizinhanças e comunidades, cuja intensidade
possa ser quantificada por ferramentas de Redes Complexas.

% ============================================================
\section*{3. Objetivo geral do projeto}

Caracterizar, por meio de redes complexas, quanto e de que forma a estrutura
de similaridade entre representações de tokens se transforma entre o espaço
lexical de entrada e as representações contextuais produzidas pelas camadas
inicial, intermediária e final de um Transformer.

% ============================================================
\section*{4. Objetivos específicos do projeto}

O projeto buscará responder às seguintes perguntas:

\begin{enumerate}[label=\textbf{P\arabic*.}, leftmargin=2.1em]

    \item \textbf{Quanto e de que forma a organização da rede muda?}
    Quantificar a reorganização local, pela similaridade de Jaccard entre as
    vizinhanças de uma mesma ocorrência e pela dominância lexical, e a
    reorganização global, pela distribuição de graus de entrada (surgimento ou
    desaparecimento de \textit{hubs}), pela clusterização, pelas distâncias e
    pelas comunidades. Isso permitirá distinguir uma mudança que apenas troca
    vizinhos locais de uma mudança que altera a organização em escala maior.

    \item \textbf{As mudanças concentram-se em determinadas camadas?}
    Comparar as transições lexical $\to$ inicial, inicial $\to$ intermediária
    e intermediária $\to$ final, identificando qual delas concentra a maior
    queda de Jaccard e de dominância lexical, e verificar se essa resposta
    depende da faixa de frequência do token e de ele ser uma palavra inteira
    ou uma subpalavra.

    \item \textbf{Surgem comunidades associadas a contextos semelhantes?}
    Verificar se as comunidades deixam de ser explicadas pelo tipo de token e
    passam a ser explicadas pelo contexto de origem (documento, sentença ou
    tema), usando medidas de concordância entre partições e inspeção
    qualitativa das comunidades.

    \item \textbf{Ocorrências de um mesmo token separam-se conforme o
    contexto?}
    Para tipos com múltiplas ocorrências, especialmente os encontrados em
    contextos distintos, verificar se representações inicialmente idênticas
    passam a ocupar comunidades ou regiões diferentes da rede.

\end{enumerate}

Para responder a essas perguntas, serão realizadas as seguintes etapas:

\begin{itemize}

    \item Construir uma \textbf{rede lexical de tipos}, uma \textbf{rede
    lexical por ocorrência} e \textbf{redes contextuais} com aproximadamente
    $1.5\times10^4$ ocorrências, mantendo os mesmos vértices em todas as redes
    por ocorrência para permitir comparação direta.

    \item Calcular e comparar as principais propriedades das redes:
    $|V|$, $|E|$, grau médio, distribuição de graus (de entrada e de saída na
    versão direcionada), densidade, clusterização global e local, distância
    média, diâmetro, componentes conexos e distribuição de seus tamanhos.

    \item Reportar os resultados por faixa de frequência e por categoria de
    token (palavra inteira ou subpalavra), avaliando sua robustez ao tratamento
    de empates, ao valor de $k$ e à regra de simetrização.

    \item Visualizar a evolução da estrutura das redes ao longo das camadas.

\end{itemize}

% ============================================================
\section*{5. Terminologia: palavras, tokens e ocorrências}

Para permitir a interpretação dos resultados, serão distinguidos:

\begin{itemize}

    \item \textbf{Palavra:} unidade do texto delimitada por espaços e
    pontuação, antes da tokenização.

    \item \textbf{Token:} unidade produzida pelo tokenizer do modelo, que
    pode corresponder a uma palavra inteira, a um fragmento de palavra
    (subpalavra) ou a um sinal de pontuação. No piloto, por exemplo, a palavra
    ``praca'' foi dividida nos tokens ``pra'' e ``ca''.

    \item \textbf{Tipo de token:} um elemento do vocabulário (\texttt{token\_id}),
    denotado por $t$. Seu número de ocorrências na amostra é $f_t$.

    \item \textbf{Ocorrência:} uma instância posicional de um tipo de token em
    um texto, denotada por $i$, com tipo $t_i$.

\end{itemize}

Cada ocorrência será rotulada como \textit{palavra inteira} (a palavra
corresponde a um único token), \textit{início de palavra fragmentada},
\textit{continuação de palavra} ou \textit{pontuação}. As interpretações
semânticas e contextuais serão feitas preferencialmente sobre tokens que são
palavras inteiras de conteúdo, pois um fragmento como ``ca'' não tem
significado próprio. As métricas serão reportadas separadamente para cada
categoria.

% ============================================================
\section*{6. Origem dos dados}

Será utilizada uma amostra de textos reais em português, inicialmente obtida da
Wikipédia. Os textos serão processados utilizando o tokenizer do
próprio Transformer analisado. Tokens especiais serão excluídos.

Para cada ocorrência serão mantidos o token, o \texttt{token\_id}, a posição,
o documento, a sentença, a palavra de origem, a posição do token dentro da
palavra e o contexto textual. Esses metadados são necessários para rotular as
ocorrências (Seção~5) e para associar comunidades a contextos (P3).

A análise principal utilizará aproximadamente $1.5\times10^4$ ocorrências de
tokens. Para impedir que tokens muito frequentes dominem artificialmente a
topologia, será imposto um limite máximo $f_{\max}$ de ocorrências por tipo.
Os tipos serão agrupados em \textbf{faixas de frequência} na amostra,
inicialmente $f_t \in \{1\text{--}2\}$ (raros), $\{3\text{--}9\}$,
$\{10\text{--}49\}$ e $\{\ge 50\}$ (até $f_{\max}$), ajustadas à distribuição
observada.

Tokens raros permanecem como vértices, para que o conjunto de vértices seja o
mesmo em todas as redes, mas serão excluídos das análises por tipo. Para
$f_t=1$ a dominância lexical é nula por definição, e para $f_t=2$ ela não
passa de $1/k$. Por isso, dominância lexical e análise intra-token serão
calculadas apenas para $f_t\ge 3$, com a análise de separação em comunidades
(P4) concentrada em $f_t\ge 10$. Como análise complementar, poderá ser
selecionado um subconjunto controlado de palavras utilizadas em contextos
distintos.

% ============================================================
\section*{7. Como os objetivos serão atingidos?}

\subsection*{7.1 Representações e redes}

Será utilizado um Transformer aberto, preferencialmente entre
4B e 8B parâmetros, com $L$ blocos. Para cada ocorrência $i$, serão extraídos
seu embedding lexical $e_{t_i}$, lido diretamente da matriz de embeddings de
entrada, e os estados internos $h_i^{(\ell)}$ na saída de três camadas
representativas: o primeiro bloco ($\ell=1$), um bloco intermediário
($\ell=\lfloor L/2\rfloor$) e o último bloco ($\ell=L$).

Serão construídos três tipos de rede:

\begin{enumerate}

    \item \textbf{Rede lexical de tipos:} um vértice para cada tipo de token,
    representado por seu embedding de entrada $e_t$ e com a frequência $f_t$
    como atributo.

    \item \textbf{Rede lexical por ocorrência:} os vértices são as ocorrências
    analisadas, mas cada ocorrência $i$ é representada por
    $x_i=e_{t_i}$. Assim, ocorrências de um mesmo tipo possuem inicialmente
    representações idênticas.

    \item \textbf{Redes contextuais:} os mesmos vértices são representados
    pelos estados $x_i=h_i^{(\ell)}$ em cada camada $\ell$ analisada.

\end{enumerate}

A similaridade entre dois vértices será medida pelo cosseno,
\[
s_{ij} =
\frac{x_i^\top x_j}
{\lVert x_i\rVert\,\lVert x_j\rVert},
\]
e cada vértice $i$ escolherá o conjunto $N_i$ de seus $k$ vizinhos mais
próximos (excluindo o próprio vértice). O mesmo $k$ e as mesmas regras serão
utilizados em todas as redes, e diferentes valores de $k$ serão testados como
análise de robustez. Busca aproximada com FAISS poderá ser utilizada para
reduzir o custo computacional.

\subsection*{7.2 Direcionalidade}

A relação de $k$ vizinhos mais próximos não é simétrica: $j\in N_i$ não
implica $i\in N_j$. O objeto produzido pelo k-NN é, portanto, uma
\textbf{rede direcionada} com arestas $i\to j$ para $j\in N_i$, na qual todo
vértice tem grau de saída exatamente $k$. Essa versão direcionada será usada
para:

\begin{itemize}
    \item as medidas de vizinhança (Jaccard e dominância lexical), que
    dependem dos conjuntos $N_i$ de tamanho fixo $k$;
    \item a distribuição de graus de entrada, que revela \textit{hubs}
    (vértices escolhidos como vizinhos por muitos outros), e a reciprocidade.
\end{itemize}

Para clusterização, distâncias, diâmetro, componentes conexos e comunidades,
será usada a \textbf{versão não direcionada por união}, na qual existe a
aresta $\{i,j\}$ se $j\in N_i$ ou $i\in N_j$; nela, o grau de cada vértice é
pelo menos $k$. A versão por k-NN mútuo ($j\in N_i$ e $i\in N_j$) será usada
como análise de robustez. Distância média e diâmetro serão calculados no maior
componente conexo.

\subsection*{7.3 Tratamento de empates}

Na rede lexical por ocorrência, as $f_{t_i}-1$ outras ocorrências do mesmo
tipo têm similaridade $s_{ij}=1$ com $i$. Se $f_{t_i}-1>k$, a escolha de quais
delas entram em $N_i$ depende de um desempate arbitrário. Se $f_{t_i}-1<k$, as
vagas restantes são ocupadas por ocorrências do tipo mais próximo, que também
são idênticas entre si e, portanto, também empatam. Como os resultados das
redes contextuais são comparados com essa rede inicial, os empates serão
tratados explicitamente de três formas:

\begin{enumerate}[label=(\alph*)]

    \item \textbf{Desempate aleatório com sementes (principal):} entre
    candidatos empatados na fronteira do $k$-ésimo vizinho, a escolha será
    uniforme ao acaso. A construção será repetida com $R$ sementes (por
    exemplo, $R=20$), e todas as métricas serão reportadas como média $\pm$
    desvio padrão. O Jaccard entre as redes lexicais de duas sementes
    distintas servirá como nível de referência: uma mudança lexical $\to$
    contextual só será interpretada como reorganização se for maior que a
    variação produzida pelo próprio desempate. O desempate pela menor posição
    no texto, usado no piloto, será mantido como caso determinístico de
    referência.

    \item \textbf{Inclusão de todos os empatados:} cada vértice se conecta a
    todos os candidatos com similaridade maior ou igual à do $k$-ésimo
    vizinho. Essa variante elimina a arbitrariedade, ao custo de um grau de
    saída variável, e será usada como análise de robustez.

    \item \textbf{Vizinhos de tipos distintos:} as ocorrências do próprio tipo
    são ignoradas, e cada ocorrência $i$ passa a ter como vizinhança o conjunto
    $T_i$ dos $k$ tipos distintos ($t\neq t_i$) mais próximos, considerando,
    para cada tipo, sua ocorrência mais similar a $i$. Na etapa lexical, $T_i$
    é simplesmente o conjunto dos $k$ tipos mais próximos de $t_i$, o que
    coincide com a rede lexical de tipos, com $f_t$ como atributo, e não possui
    empates. Nas camadas contextuais, $T_i$ indica \emph{de quais outros
    tokens} a ocorrência se aproxima em seu contexto. Essa variante
    complementa a dominância lexical, que continua sendo calculada nas
    variantes (a) e (b).

\end{enumerate}

Nas camadas contextuais, empates exatos são praticamente inexistentes;
similaridades que diferirem por menos de uma tolerância numérica $\varepsilon$
(por exemplo, $10^{-6}$) serão tratadas como empates, pela mesma regra.

\subsection*{7.4 Medidas}

A reorganização da vizinhança de uma ocorrência entre as redes $a$ e $b$ será
medida pela similaridade de Jaccard
\[
J_i(a,b)=\frac{|N_i^{(a)}\cap N_i^{(b)}|}{|N_i^{(a)}\cup N_i^{(b)}|},
\]
calculada também sobre os conjuntos de tipos $T_i$ na variante (c). A
dominância lexical será a fração de vizinhos do mesmo tipo,
$D_i^{(\ell)}=|\{j\in N_i^{(\ell)}: t_j=t_i\}|/k$. Comunidades serão
identificadas com Leiden ou Louvain na rede não direcionada e comparadas entre
camadas por NMI e ARI. Para responder a P3, as partições também serão
comparadas com as partições por tipo de token e por documento/sentença, e
será medida a pureza de cada comunidade em relação a esses rótulos.

% ============================================================
\section*{8. Como os resultados serão avaliados?}

O cumprimento dos objetivos será avaliado em três níveis. Primeiro, será
verificado se todas as redes propostas puderam ser construídas com escala
superior a $10^4$ vértices e se as métricas estruturais requeridas pela
disciplina foram calculadas.

Segundo, as perguntas P1--P3 serão respondidas comparando as redes entre
camadas. Para P1, a distribuição de $J_i$ e de $D_i$ indicará quanto a
identidade lexical continua determinando a vizinhança, e as métricas globais e
a distribuição de graus de entrada mostrarão se a transformação também ocorre
em escala maior. Para P2, a transição com a maior queda de Jaccard e de
dominância indicará onde a reorganização se concentra. Para P3, a organização
será considerada contextual se a concordância das comunidades com o tipo de
token diminuir e a concordância com documento/sentença aumentar ao longo das
camadas.

Por fim, para P4, serão analisados tipos com múltiplas ocorrências
($f_t\ge 10$), verificando em quantas comunidades suas ocorrências se
distribuem e como a similaridade média entre elas evolui.

Todos os resultados serão reportados por faixa de frequência e por categoria
de token, como média $\pm$ desvio entre as $R$ sementes de desempate, e
repetidos para diferentes valores de $k$, para as três variantes de tratamento
de empates e para as simetrizações por união e mútua. Isso reduz a
dependência das conclusões de uma única parametrização.

% ============================================================
\section*{9. Planejamento e cronograma}

\begin{itemize}

    \item \textbf{Semanas 1--2:} seleção do modelo e corpus, definição da
    amostragem e das faixas de frequência, rotulagem palavra/subpalavra e
    extração dos embeddings e estados internos.

    \item \textbf{Semanas 2--3:} construção das redes direcionadas e não
    direcionadas, com as três variantes de tratamento de empates e as $R$
    sementes.

    \item \textbf{Semanas 3--4:} cálculo das métricas estruturais, persistência
    de vizinhança e dominância lexical, por faixa de frequência e categoria.

    \item \textbf{Semanas 4--5:} detecção de comunidades, associação com
    contextos, análise intra-token, visualizações e testes de robustez.

    \item \textbf{Semanas 5--6:} consolidação dos resultados e redação do
    relatório final.

\end{itemize}

% ============================================================
\section*{10. Trabalhos futuros}

\begin{itemize}

    \item Estender a análise para todas as camadas do modelo, comparando cada
    camada com a seguinte ($\ell\to\ell+1$), para localizar com mais precisão
    onde a reorganização ocorre.

    \item Construir redes no nível de palavras, agregando as representações
    das subpalavras que compõem cada palavra.

\end{itemize}

% ============================================================
\section*{11. Resposta às observações da avaliação}

\begin{center}
\small
\begin{tabularx}{\textwidth}{@{}lX@{}}
\toprule
\textbf{Observação} & \textbf{Onde foi tratada} \\
\midrule
Objetivos além de ``medir'' e ``comparar'' &
Seção 4: perguntas P1--P4 e forma de respondê-las na Seção 8. \\
(1) Rede direcionada ou não &
Seção 7.2: k-NN direcionado para vizinhanças e graus de entrada;
não direcionado por união para as demais métricas; k-NN mútuo como robustez. \\
(2) Faixas de frequência e tokens raros &
Seção 6: faixas de frequência, limite $f_{\max}$ e exclusão de
tokens com $f_t\le 2$ das análises por tipo. \\
(3) Empates na rede inicial &
Seção 7.3: desempate aleatório com $R$ sementes, inclusão de empatados e
vizinhos de tipos distintos. \\
(4) Palavras e tokens &
Seção 5: definições e rotulagem de cada ocorrência como palavra inteira ou
subpalavra. \\
\bottomrule
\end{tabularx}
\end{center}

\end{document}
